% GENERATED FILE. Edit canonical lessons, then run npm run generate:books.
% canonical-source-hash: fnv1a64:96dbf680c70af296
% canonical-lessons: SA-C367-sambhavah, SA-C367-sambhavayati, SA-C367-dhruvam, SA-C367-ksipram, SA-C367-apeksate

\chapter{Is It Possible?}
\label{ch:sa-is-it-possible}
\hlchaptermodality{367}

\begin{chapteropening}
\textbf{By the end of this chapter:} \emph{I can say whether something is possible and when it may happen.}

Complete the last lesson of chapter 367: I can say whether something is possible and when it may happen.
\end{chapteropening}

\section[\texorpdfstring{\sk{सम्भवः} (sambhavaḥ)}{सम्भवः (sambhavaḥ)}]{\sk{सम्भवः} (sambhavaḥ) — a possibility}

\label{lesson:SA-C367-sambhavah}



\begin{quote}
Before the new one: say the Sanskrit for considers, then the Sanskrit for ascertains.
\end{quote}

\subsection*{You'll want to know: \sk{सम्भवः}}
\textbf{\sk{सम्भवः}} — \emph{sambhavaḥ} — \textquotedblleft{}a possibility\textquotedblright{}.

\begin{quote}
\sk{यदि} \sk{धनम्} \sk{अस्ति}, \sk{तर्हि} \sk{यात्रायाः} \sk{सम्भवः} \sk{अस्ति।}
\end{quote}

\begin{quote}
\emph{yadi dhanam asti, tarhi yātrāyāḥ sambhavaḥ asti.}
\end{quote}

\begin{quote}
If there is money, a journey is possible.
\end{quote}

A masculine noun.

\begin{cousinweb}
\textbf{\sk{सम्भवः}} (\emph{sambhavaḥ}) is built on \emph{sam}, \textquotedblleft{}together\textquotedblright{}, and the root √bhū, \textquotedblleft{}to be\textquotedblright{}; MW gives \emph{sam-bhava} \textquotedblleft{}capacity, ability, possibility\textquotedblright{} (MW).
\end{cousinweb}

\begin{grammarlens}[title={What you've built}]
One more word for saying what happens if.
\end{grammarlens}

\subsection*{Your turn}
\begin{itemize}
\raggedright
  \item \emph{Say it:} \emph{sambhavaḥ}
  \item \emph{Say it:} \emph{sambhavaḥ}, once more
  \item \emph{Say it:} \emph{yadi dhanam asti, tarhi yātrāyāḥ sambhavaḥ asti.} — If there is money, a journey is possible.
  \item \emph{Recall:} say the Sanskrit for considers, then the Sanskrit for ascertains, then say \emph{sambhavaḥ} again
\end{itemize}

\begin{tcolorbox}[breakable,colback=teal!4,colframe=teal!35!black,title={Before you move on}]
What does \sk{सम्भवः} mean? (\textquotedblleft{}A possibility\textquotedblright{}.) Say it once more.
\end{tcolorbox}
\section[\texorpdfstring{\sk{सम्भावयति} (sambhāvayati)}{सम्भावयति (sambhāvayati)}]{\sk{सम्भावयति} (sambhāvayati) — thinks likely; supposes}

\label{lesson:SA-C367-sambhavayati}



\begin{quote}
Before the new one: say the Sanskrit for ascertains, then the Sanskrit for a possibility.
\end{quote}

\subsection*{You'll want to know: \sk{सम्भावयति}}
\textbf{\sk{सम्भावयति}} — \emph{sambhāvayati} — \textquotedblleft{}thinks likely; supposes\textquotedblright{}.

\begin{quote}
\sk{यदि} \sk{आकाशः} \sk{कृष्णः}, \sk{तर्हि} \sk{अहं} \sk{वर्षां} \sk{सम्भावयामि।}
\end{quote}

\begin{quote}
\emph{yadi ākāśaḥ kṛṣṇaḥ, tarhi ahaṃ varṣāṃ sambhāvayāmi.}
\end{quote}

\begin{quote}
If the sky is dark, I think rain likely.
\end{quote}

Past (he …): \sk{सः} \sk{वर्षां} \sk{सम्भावितवान्}, \emph{saḥ varṣāṃ sambhāvitavān}.

\begin{cousinweb}
\textbf{\sk{सम्भावयति}} (\emph{sambhāvayati}) is built on \emph{sam} and the root √bhū, \textquotedblleft{}to be\textquotedblright{}, in its causative form; MW gives the senses \textquotedblleft{}to suggest a possibility\textquotedblright{} and \textquotedblleft{}to suppose anything possible\textquotedblright{} (MW).
\end{cousinweb}

\begin{grammarlens}[title={What you've built}]
One more word for saying what happens if.
\end{grammarlens}

\subsection*{Your turn}
\begin{itemize}
\raggedright
  \item \emph{Say it:} \emph{sambhāvayati}
  \item \emph{Say it:} \emph{sambhāvayati}, once more
  \item \emph{Say it:} \emph{yadi ākāśaḥ kṛṣṇaḥ, tarhi ahaṃ varṣāṃ sambhāvayāmi.} — If the sky is dark, I think rain likely.
  \item \emph{Recall:} say the Sanskrit for ascertains, then the Sanskrit for a possibility, then say \emph{sambhāvayati} again
\end{itemize}

\begin{tcolorbox}[breakable,colback=teal!4,colframe=teal!35!black,title={Before you move on}]
What does \sk{सम्भावयति} mean? (\textquotedblleft{}Thinks likely; supposes\textquotedblright{}.) Say it once more.
\end{tcolorbox}
\section[\texorpdfstring{\sk{ध्रुवम्} (dhruvam)}{ध्रुवम् (dhruvam)}]{\sk{ध्रुवम्} (dhruvam) — for certain; surely}

\label{lesson:SA-C367-dhruvam}



\begin{quote}
Before the new one: say the Sanskrit for a possibility, then the Sanskrit for thinks likely.
\end{quote}

\subsection*{You'll want to know: \sk{ध्रुवम्}}
\textbf{\sk{ध्रुवम्}} — \emph{dhruvam} — \textquotedblleft{}for certain; surely\textquotedblright{}.

\begin{quote}
\sk{यदि} \sk{भवान्} \sk{आगच्छति}, \sk{तर्हि} \sk{सः} \sk{ध्रुवं} \sk{प्रसन्नः} \sk{भविष्यति।}
\end{quote}

\begin{quote}
\emph{yadi bhavān āgacchati, tarhi saḥ dhruvaṃ prasannaḥ bhaviṣyati.}
\end{quote}

\begin{quote}
If you come, he will certainly be pleased.
\end{quote}

\begin{cousinweb}
\textbf{\sk{ध्रुवम्}} (\emph{dhruvam}) is built on \emph{dhruva}, \textquotedblleft{}fixed, firm\textquotedblright{}; MW gives \emph{dhruvam} \textquotedblleft{}firmly, constantly, certainly, surely\textquotedblright{} (MW).
\end{cousinweb}

\begin{grammarlens}[title={What you've built}]
One more word for saying what happens if.
\end{grammarlens}

\subsection*{Your turn}
\begin{itemize}
\raggedright
  \item \emph{Say it:} \emph{dhruvam}
  \item \emph{Say it:} \emph{dhruvam}, once more
  \item \emph{Say it:} \emph{yadi bhavān āgacchati, tarhi saḥ dhruvaṃ prasannaḥ bhaviṣyati.} — If you come, he will certainly be pleased.
  \item \emph{Recall:} say the Sanskrit for a possibility, then the Sanskrit for thinks likely, then say \emph{dhruvam} again
\end{itemize}

\begin{tcolorbox}[breakable,colback=teal!4,colframe=teal!35!black,title={Before you move on}]
What does \sk{ध्रुवम्} mean? (\textquotedblleft{}For certain; surely\textquotedblright{}.) Say it once more.
\end{tcolorbox}
\section[\texorpdfstring{\sk{क्षिप्रम्} (kṣipram)}{क्षिप्रम् (kṣipram)}]{\sk{क्षिप्रम्} (kṣipram) — swiftly; without delay}

\label{lesson:SA-C367-ksipram}



\begin{quote}
Before the new one: say the Sanskrit for thinks likely, then the Sanskrit for for certain.
\end{quote}

\subsection*{You'll want to know: \sk{क्षिप्रम्}}
\textbf{\sk{क्षिप्रम्}} — \emph{kṣipram} — \textquotedblleft{}swiftly; without delay\textquotedblright{}.

\begin{quote}
\sk{यदि} \sk{वायुः} \sk{वहति}, \sk{तर्हि} \sk{क्षिप्रं} \sk{वर्षा} \sk{भविष्यति।}
\end{quote}

\begin{quote}
\emph{yadi vāyuḥ vahati, tarhi kṣipraṃ varṣā bhaviṣyati.}
\end{quote}

\begin{quote}
If the wind blows, rain will come swiftly.
\end{quote}

\begin{cousinweb}
\textbf{\sk{क्षिप्रम्}} (\emph{kṣipram}) comes from the root √kṣip, \textquotedblleft{}to throw\textquotedblright{} (CDIAL 3688 sets \emph{kṣiprá} under it); MW gives \emph{kṣipra} \textquotedblleft{}springing\textquotedblright{} and \textquotedblleft{}quick, speedy, swift\textquotedblright{}, and \emph{kṣipram} \textquotedblleft{}quickly, immediately\textquotedblright{} (MW).
\end{cousinweb}

\begin{grammarlens}[title={What you've built}]
One more word for saying what happens if.
\end{grammarlens}

\subsection*{Your turn}
\begin{itemize}
\raggedright
  \item \emph{Say it:} \emph{kṣipram}
  \item \emph{Say it:} \emph{kṣipram}, once more
  \item \emph{Say it:} \emph{yadi vāyuḥ vahati, tarhi kṣipraṃ varṣā bhaviṣyati.} — If the wind blows, rain will come swiftly.
  \item \emph{Recall:} say the Sanskrit for thinks likely, then the Sanskrit for for certain, then say \emph{kṣipram} again
\end{itemize}

\begin{tcolorbox}[breakable,colback=teal!4,colframe=teal!35!black,title={Before you move on}]
What does \sk{क्षिप्रम्} mean? (\textquotedblleft{}Swiftly; without delay\textquotedblright{}.) Say it once more.
\end{tcolorbox}
\section[\texorpdfstring{\sk{अपेक्षते} (apekṣate)}{अपेक्षते (apekṣate)}]{\sk{अपेक्षते} (apekṣate) — expects}

\label{lesson:SA-C367-apeksate}



\begin{quote}
Before the new one: say the Sanskrit for for certain, then the Sanskrit for swiftly.
\end{quote}

\subsection*{You'll want to know: \sk{अपेक्षते}}
\textbf{\sk{अपेक्षते}} — \emph{apekṣate} — \textquotedblleft{}expects\textquotedblright{}.

\begin{quote}
\sk{यदि} \sk{वर्षा} \sk{भवति}, \sk{तर्हि} \sk{कृषकः} \sk{उत्तमं} \sk{धान्यम्} \sk{अपेक्षते।}
\end{quote}

\begin{quote}
\emph{yadi varṣā bhavati, tarhi kṛṣakaḥ uttamaṃ dhānyam apekṣate.}
\end{quote}

\begin{quote}
If it rains, the farmer expects a good harvest.
\end{quote}

Past (he …): \sk{सः} \sk{उत्तरम्} \sk{अपेक्षितवान्}, \emph{saḥ uttaram apekṣitavān}.

\begin{cousinweb}
\textbf{\sk{अपेक्षते}} (\emph{apekṣate}) is built on \emph{apa}, \textquotedblleft{}away\textquotedblright{}, and the root √īkṣ, \textquotedblleft{}to see\textquotedblright{}; MW gives \emph{apekṣ} \textquotedblleft{}to look for, to wait for, to expect\textquotedblright{} (MW).
\end{cousinweb}

\begin{grammarlens}[title={What you've built}]
One more word for saying what happens if.
\end{grammarlens}

\subsection*{Your turn}
\begin{itemize}
\raggedright
  \item \emph{Say it:} \emph{apekṣate}
  \item \emph{Say it:} \emph{apekṣate}, once more
  \item \emph{Say it:} \emph{yadi varṣā bhavati, tarhi kṛṣakaḥ uttamaṃ dhānyam apekṣate.} — If it rains, the farmer expects a good harvest.
  \item \emph{Recall:} say the Sanskrit for for certain, then the Sanskrit for swiftly, then say \emph{apekṣate} again
\end{itemize}

\begin{tcolorbox}[breakable,colback=teal!4,colframe=teal!35!black,title={Before you move on}]
What does \sk{अपेक्षते} mean? (\textquotedblleft{}Expects\textquotedblright{}.) Say it once more.
\end{tcolorbox}
